\section{Empirical Strategy}

The empirical design is organized as a six-layer evidence chain.

Table 1 replicates the Wang-style benchmark by testing whether high-markup sectors receive more preferential credit. Table 2 reports the distribution and stability of downstream input-demand wedges. Tables 3A--3C validate whether the wedge behaves like buyer power by relating it to downstream concentration, buyer size, supplier fragmentation, and payment pressure.

Table 4 tests whether upstream firms exposed to high-wedge downstream industries face greater cash-flow, payment, financing, and investment pressure. Table 5 tests policy targeting: whether preferential credit or policy support flows to high-exposure upstream sectors after controlling for their own markup. Table 6A tests policy effectiveness by examining whether credit support improves financial and real outcomes for exposed upstream firms. Table 6B tests the trade-off by examining whether support also increases payment persistence, buyer dependence, or weak-firm preservation.

The interpretation is intentionally sequential. Table 5 can show allocation, but it cannot establish effectiveness. Policy effectiveness requires Table 6A, and a Wang-style trade-off requires Table 6B.

The main specifications also separate downstream-wedge exposure from three competing mechanisms. First, to distinguish the paper from China's vertical-structure mechanism, Tables 4--6 control for upstreamness, upstream SOE share, upstream concentration, SOE concentration, own markup, and upstream profitability. Second, to avoid confusing exposure targeting with upstream SOE policy spillovers, the policy-targeting specifications include upstream SOE policy-support controls whenever subsidy, tax-preference, or interest-advantage measures are available. Third, to rule out GVC upgrading and firm lifecycle dynamics, firm-level specifications control for productivity, size, age, export exposure, and processing-trade or production-line-position proxies when data allow.

Robustness checks compare the explanatory power of downstream-wedge exposure with generic network centrality, including upstreamness, Leontief influence, sales centrality, and input centrality. The key coefficient is informative only if exposure remains predictive after these vertical-structure, SOE-policy, GVC-position, and network-centrality controls.
